\chapter{Cơ sở lý thuyết và công nghệ}
\label{ch:coso}

\begin{muctieu}
Chương này tóm lược những khái niệm trực tiếp được sử dụng trong hệ thống:
phân loại ảnh, Transformer, Vision Transformer, các độ đo đánh giá và nền
tảng phần mềm triển khai.
\end{muctieu}

\section{Bài toán phân loại ảnh nhiều lớp}

Với ảnh đầu vào $x$ và tập nhãn hữu hạn $\mathcal{Y}=\{1,\ldots,K\}$, bộ
phân loại học hàm $f:x\mapsto y$ để chọn một nhãn $y\in\mathcal{Y}$. Trong
đồ án, $K=38$ và mỗi nhãn biểu diễn một món ăn. Model sinh vector điểm thô
$\mathbf{z}\in\mathbb{R}^{K}$. Xác suất lớp $k$ được tính bằng Softmax:
\begin{equation}
p(y=k\mid x)=\frac{\exp(z_k)}{\sum_{j=1}^{K}\exp(z_j)}.
\label{eq:softmax}
\end{equation}
Nhãn có xác suất lớn nhất là Top~1; ba nhãn lớn nhất tạo thành Top~3. Xác
suất cao không đồng nghĩa chắc chắn tuyệt đối, vì model vẫn có thể tự tin ở
ảnh ngoài phân phối hoặc các món có hình thức gần nhau.

\section{Transformer và cơ chế tự chú ý}

Transformer thay thế xử lý tuần tự bằng cơ chế self-attention, cho phép mỗi
token tổng hợp thông tin từ các token khác \cite{vaswani2017attention}. Với
ma trận truy vấn $Q$, khóa $K$ và giá trị $V$, scaled dot-product attention
được tính theo:
\begin{equation}
\operatorname{Attention}(Q,K,V)=
\operatorname{softmax}\left(\frac{QK^{\mathsf T}}{\sqrt{d_k}}\right)V,
\label{eq:attention}
\end{equation}
trong đó $d_k$ là số chiều của vector khóa. Multi-head attention thực hiện
nhiều phép attention song song để học các quan hệ khác nhau.

\section{Vision Transformer}

Vision Transformer (ViT) chia ảnh thành các patch không chồng lấn, biến mỗi
patch thành một token, thêm positional embedding và token phân lớp
\cite{dosovitskiy2021vit}. Ảnh 384$\times$384 với patch 16$\times$16 tạo
ra $24\times24=576$ patch. Chuỗi 576 token ảnh cùng token phân lớp đi qua 24
Transformer encoder của cấu hình ViT-L/16. Vector ở token phân lớp được đưa
vào lớp tuyến tính để sinh logits cho 38 món.

ViT tận dụng tốt học chuyển giao khi được tiền huấn luyện trên dữ liệu lớn.
Đồ án dùng model \code{VietTung04/vit-l-16-food-classifier}, một phiên bản
ViT-L/16 đã được fine-tune cho món Việt và công bố theo giấy phép Apache-2.0
\cite{vietTungModel}. Model được tải một lần, lưu local và chạy ở chế độ
\code{eval} kết hợp \code{torch.inference\_mode}.

\section{Các độ đo đánh giá}

Với mỗi lớp, $TP$, $FP$ và $FN$ lần lượt là số dự đoán đúng, dương tính giả
và âm tính giả. Các độ đo được tính theo:
\begin{align}
\mathrm{Accuracy} &= \frac{\text{số dự đoán đúng}}{\text{tổng số mẫu}},\\
\mathrm{Precision} &= \frac{TP}{TP+FP},\\
\mathrm{Recall} &= \frac{TP}{TP+FN},\\
F_1 &= \frac{2\,\mathrm{Precision}\,\mathrm{Recall}}
{\mathrm{Precision}+\mathrm{Recall}}.
\end{align}
Macro average là trung bình không trọng số của 30 lớp đánh giá, nhờ đó mỗi
lớp có đóng góp ngang nhau dù số ảnh khác nhau. Ma trận nhầm lẫn cho biết các
cặp lớp model thường dự đoán lẫn nhau.

\section{Kiến trúc client--server và REST API}

Ứng dụng Flutter đóng vai trò client, còn FastAPI cung cấp tài nguyên qua
HTTP/JSON và nhận ảnh dưới dạng multipart. FastAPI dựa trên chuẩn OpenAPI và
tự sinh Swagger UI, thuận tiện cho kiểm thử contract \cite{fastapiDocs}.
Flutter khuyến nghị tách UI, trạng thái, repository và service để giảm phụ
thuộc và tăng khả năng kiểm thử \cite{flutterArchitecture}. Dự án hiện dùng
screen/widget ở lớp UI, \code{AppState} cho trạng thái và \code{ApiClient}
cùng \code{AiService} cho lớp dữ liệu.

\section{Lưu trữ và bảo mật}

SQLite là database nhúng, serverless, giao dịch và lưu toàn bộ dữ liệu trong
một tệp \cite{sqliteAbout}. SQLAlchemy ánh xạ lớp Python sang bảng quan hệ,
quản lý session, transaction và relationship \cite{sqlalchemyDocs}. Mật khẩu
được băm bằng Argon2, một hàm memory-hard được mô tả trong RFC~9106
\cite{rfc9106}. Access token và refresh token tuân theo cấu trúc JWT của
RFC~7519 \cite{rfc7519}; refresh token còn có phiên trong database để hỗ trợ
xoay vòng và thu hồi.

\section{Công cụ và thư viện}

\begin{table}[htbp]
\centering
\caption{Các thành phần chính của hệ thống.}
\label{tab:congcu}
\small
\begin{tabularx}{\textwidth}{>{\bfseries\color{daunavy}}l >{\color{daumaroon}}l X}
\toprule
\textbf{Thành phần} & \textbf{Phiên bản kiểm chứng} & \textbf{Vai trò} \\
\midrule
Flutter & 3.47.2 & Ứng dụng mobile đa nền tảng \\
Dart & 3.13.x & Ngôn ngữ của Flutter \\
Python & 3.13.7 & Backend, audit và suy luận AI \\
FastAPI & 0.141.1 & REST API, validation và OpenAPI \\
PyTorch & 2.14.0+cpu & Tensor và suy luận model \cite{paszke2019pytorch} \\
Transformers & 5.17.0 & Processor và ViT \cite{huggingfaceTransformers} \\
SQLAlchemy & 2.1.3 & ORM và transaction \\
SQLite & 3.x & Database cục bộ \\
Android Studio & Bản cài hiện tại & Emulator, build và chạy Flutter \\
\bottomrule
\end{tabularx}
\end{table}

Phiên bản trong Bảng~\ref{tab:congcu} mô tả môi trường đã kiểm chứng của đồ
án; giới hạn dependency chính xác vẫn được quản lý trong
\code{requirements.txt} và \code{pubspec.yaml}.

\section{Dữ liệu và các công trình liên quan}

Food-101 đặt nền móng cho đánh giá nhận diện món ăn quy mô lớn với 101.000
ảnh thuộc 101 lớp \cite{bossard2014food101}. 30VNFoods tập trung vào 30 món
Việt và công bố 25.136 ảnh \cite{do2021_30vnfoods}. Model sử dụng trong đồ
án kế thừa ViT-L/16 và được tác giả model fine-tune trên 30VNFoods kết hợp dữ
liệu tự thu thập \cite{vietTungModel}. Khác với công trình chỉ tập trung phân
loại, Vị Việt tích hợp model vào luồng ứng dụng có xác thực, lịch sử, phản
hồi và kho công thức có metadata nguồn.

Với dữ liệu công thức, tám món được dịch và chuẩn hóa từ Wikibooks theo điều
kiện CC BY-SA 4.0 \cite{wikibooksCopyright}. Ảnh giao diện được chọn theo
metadata riêng của từng tệp trên Wikimedia Commons/Openverse; việc sử dụng
ảnh tuân theo yêu cầu attribution và giấy phép của tệp \cite{commonsReuse}.
